%! Dividing Polynomials; Remainder and Factor Theorems — Unit 3, Lesson 3.3 — Student Packet Cover Sheet
\documentclass[10pt]{article}
\usepackage{algebra2-article}
\usepackage{algebra2-boxes}
\usepackage{ltablex}
\keepXColumns

\begin{document}

% Full-bleed forest banner
\begin{tikzpicture}[remember picture, overlay]
  \fill[forest] (current page.north west)
    rectangle ([yshift=-0.9in]current page.north east);
\end{tikzpicture}
\vspace{-0.8in}

\begin{center}
  {\color{white}\textbf{\LARGE Algebra 2}}\\[4pt]
  {\color{forestmid}\large\textbf{Unit 3 \quad Polynomial Functions}}\\[6pt]
  {\color{white}\textbf{\large Lesson 3.3 \quad Dividing Polynomials; Remainder \& Factor Theorems}}
\end{center}

\vspace{0.22in}
\namedateperiod
\vspace{0.18in}

\begin{learningtargetbox}
\small
\begin{enumerate}[leftmargin=1.5em, itemsep=5pt]
  \item \ldots{} divide a polynomial by a \textbf{monomial} by splitting the fraction term by term, in one and two variables.
  \item \ldots{} use \textbf{long division} to divide by a binomial or a factorable trinomial, using a \textbf{placeholder $0$} for every missing power.
  \item \ldots{} use \textbf{synthetic division} when the divisor is $x - r$, and remember that $x+2$ means $r = -2$.
  \item \ldots{} write a division result as \textbf{quotient $+\;\frac{\text{remainder}}{\text{divisor}}$}, and check it with dividend $=$ divisor $\times$ quotient $+$ remainder.
  \item \ldots{} use the \textbf{Remainder Theorem} to find a remainder --- or evaluate $f(r)$ --- without dividing.
  \item \ldots{} use the \textbf{Factor Theorem} to decide whether $(x-r)$ is a factor, then divide it out and \textbf{finish} factoring.
\end{enumerate}
\end{learningtargetbox}

\vspace{0.14in}

\begin{tocbox}
\small
\renewcommand{\arraystretch}{1.5}
\begin{tabularx}{\linewidth}{@{} c l X r @{}}
\rowcolor{forestbg}
\textbf{\#} & \textbf{Component} & \textbf{Description} & \textbf{Score} \\
\midrule
1 & Warm-Up        & Spiral review: long division of numbers, monomial quotients, evaluating $f$ & \blank{1.2cm} \\
2 & Guided Notes   & Long division, synthetic division, and the Remainder \& Factor Theorems     & \blank{1.2cm} \\
3 & Group Activity & \emph{Divide and Conquer} --- dividing and testing factors, by tier         & \blank{1.2cm} \\
4 & Exit Ticket    & Quick check: one synthetic division, one theorem, one SOL-style item        & \blank{1.2cm} \\
5 & Homework       & Mixed division practice, factor-hunting, and a volume model                 & \blank{1.2cm} \\
\midrule
  & \multicolumn{2}{r}{\textbf{Total}} & \blank{1.2cm} \\
\end{tabularx}
\end{tocbox}

\vspace{0.10in}

\begin{remindbox}
\small Everything today rests on one sentence you have known since elementary school:
\textbf{dividend $=$ divisor $\times$ quotient $+$ remainder}. Just as $247 = 5 \cdot 49 + 2$,
every polynomial division gives $P(x) = D(x)\,Q(x) + R(x)$. Three habits keep the work clean:
write the dividend in \textbf{standard form with a placeholder $0$} for every missing power;
\textbf{subtract} (change every sign) at each long-division step; and when using synthetic division
remember that the divisor $x + 2$ means $r = \mathbf{-2}$. The payoff is the remainder. A remainder
of $0$ means the divisor is a \textbf{factor} --- and that single fact is what finally cracks the
polynomials Lesson 3.2 could not touch.
\end{remindbox}

\end{document}
